\par\medskip
\begin{center}\hypertarget{native-166-article119}{}\pdfbookmark[2]{119.}{article-119}119.\end{center}
It would be easy to use the connection between the data and the unknowns of our problem to reduce its formulation to six equations, or even fewer, since one or another unknown could be eliminated quite conveniently. But the connection is so complicated that these equations would become extremely intractable. A separation of the unknowns that ultimately yields an equation containing just one unknown may, generally speaking\footnote{Whenever the observations are so close together that the time intervals may be treated as infinitesimal quantities, such a separation does succeed, and the whole problem reduces to solving an algebraic equation of the seventh or eighth degree.}, be regarded as impossible. Still less can the complete solution of the problem be achieved by direct operations alone.

We can, however, reduce our problem in various ways to solving \emph{two} equations $X=0$, $Y=0$, in which only two unknowns $x$, $y$ remain coupled. The unknowns $x$, $y$ need not be two of the elements themselves. They may be quantities related to the elements in any way, provided that the elements can readily be derived from them once they have been found. It is also clearly unnecessary to express $X$, $Y$ as explicit functions of $x$, $y$. It suffices for a system of equations to relate them in such a way that we can pass from given values of $x$, $y$ to the corresponding values of $X$, $Y$.
